\section{Introduction}

Most machine learning datasets are invisible --- not because they lack quality, but because they lack tags. On OpenML, a platform hosting over 5,000 actively-studied datasets, we find that datasets with at least one keyword tag register 22.6\% more ML task experiments than their untagged counterparts, even after controlling for dataset size, age, and the decade of upload. This finding is not marginal: the incidence rate ratio of 1.2263 (95\% CI [1.1681, 1.2873], $p<0.001$, N=5,217) is statistically overwhelming, robust to seven model variants, and reflects a structural platform property rather than a temporal artifact.

The disparity in ML dataset adoption is a well-documented but poorly-understood phenomenon. Some datasets accumulate hundreds of task registrations and experimental runs; others, of comparable quality and scope, sit largely unused. Most explanations focus on intrinsic dataset properties --- the number of instances, dimensionality, task difficulty --- as the primary drivers of engagement. These structural features matter, and we control for them explicitly. But they cannot explain adoption gaps among datasets of similar size and complexity. Something else is at work.

The mechanism we identify is discoverability via search-indexed metadata. On platforms like OpenML, researchers discover datasets primarily through keyword search. A dataset without keyword tags is structurally absent from tag-indexed search results, regardless of its intrinsic merit. This is the FAIR F1 (Findability) principle operationalized: keyword tags are the mechanism by which datasets enter researcher awareness. The FAIR Data Principles \citep{Wilkinson2016FAIR} established this as a foundational requirement for dataset reuse, but the principle has remained largely aspirational --- stated normatively without empirical IRR quantification. We provide that quantification.

Our key insight is that the FAIR F1 tagging mechanism on OpenML operates as a \textbf{threshold-plus-amplification} structure. Binary tag presence (\texttt{has\_tags=1} vs.\ \texttt{has\_tags=0}) captures the critical search-index-membership threshold: any tagging creates a 22.6\% adoption advantage. Above this threshold, tag count magnitude amplifies discovery probability log-linearly: within the tagged subset (N=2,625), each unit increase in $\log(\text{tag\_count}+1)$ is associated with 53.3\% more task registrations (IRR=1.5332, $p<0.001$). The categorical shape of this amplification is non-uniform --- the 6+ tag tier drives a meaningfully distinct effect (IRR=1.2861) while 1--5 tags produce statistically indistinguishable effects --- suggesting that comprehensive tagging, not minimal tagging, is the behavior that maximizes discovery.

Crucially, this effect is structural rather than temporal. The mechanism survives decade fixed effects with only 12.2\% attenuation (attenuation\_ratio=1.1219), despite extreme decade-tag collinearity (Cram\'er's V=0.823). Within each upload cohort, tagged datasets attract more tasks than untagged ones --- confirming that tag-indexed search membership operates as a stable platform property across eras.

We make the following contributions:

\textbf{C1. First empirical NB-2 quantification of FAIR F1 keyword tagging $\rightarrow$ ML dataset adoption.} We provide incidence rate ratio estimates for \texttt{has\_tags} binary (IRR=1.2263) and log tag count (IRR=1.5332) in negative binomial regression appropriate for overdispersed count outcomes, on the OpenML platform with N=5,217 datasets.

\textbf{C2. Threshold-plus-amplification characterization of the FAIR F1 mechanism.} Beyond establishing that tagging predicts adoption, we characterize the mechanism shape: binary threshold (any tags $\rightarrow$ search index membership), continuous log-linear amplification (more tags $\rightarrow$ more search pathways), and non-uniform categorical tiers (6+ tag dominance). This mechanistic precision exceeds prior framework-level descriptions.

\textbf{C3. Structural mechanism verification via decade fixed effects.} We test the critical confound --- that tagged datasets are simply older, from a different platform era, and thus trivially more adopted. The \texttt{has\_tags} effect survives \texttt{C(decade)} FE with 12.2\% attenuation (vs.\ 98.6\% attenuation for our prior composite metadata score attempt on the same corpus), establishing the within-decade, structural character of the finding.

\textbf{C4. Informative negative for uniform categorical dose-response.} The H-M3 hypothesis test reveals that OpenML tagging behavior is bimodal --- users tag comprehensively (3+ tags) or not at all --- making the 1--2 tag range a sparse middle ground (N=73). This informative negative constrains future tagging intervention designs.

We organize the paper as follows: Section~\ref{sec:related} situates our contribution within the FAIR principles literature and ML metadata adoption research. Section~\ref{sec:methodology} describes the study design, corpus, and NB-2 regression methodology. Section~\ref{sec:experimental} presents our experimental approach for testing each prediction. Section~\ref{sec:results} reports results for all four hypotheses. Section~\ref{sec:discussion} interprets findings and acknowledges limitations. Section~\ref{sec:conclusion} concludes with implications and future directions.
